% !TEX root = task-1-15-5.tex
% Editable task bank. Append new BankTask blocks at the end of the document.
\documentclass[12pt,a4paper]{article}
\usepackage{physics-textbook}
\usepackage{graphicx,needspace}
\newif\ifShowSolutions
\ifdefined\TutorEdition
  \ShowSolutionstrue
\else
  \ShowSolutionsfalse % Change to \ShowSolutionstrue to display solutions.
\fi
\newcounter{banktask}
\newcommand{\BankTask}[3]{%
  \refstepcounter{banktask}%
  \Needspace{35mm}\par\medskip\noindent\begin{minipage}{\linewidth}%
  \noindent\textbf{\thebanktask.}\quad #2\par
  \ifShowSolutions
    #3\par
    \ifdefined\TutorEdition
      \ifcsname TutorSource#1\endcsname
        \smallskip{\small\csname TutorSource#1\endcsname\par}%
      \fi
    \fi
  \fi
  \end{minipage}\par\medskip
}
\begin{document}
\BookSubsection{1.15.5}{Потенциальная энергия тела вблизи поверхности Земли}


% Source sdamgia-528: https://phys-ege.sdamgia.ru/problem?id=528
\BankTask{sdamgia-528}{
Камень массой $1\,\text{кг}$ брошен вертикально вверх. В начальный момент его энергия равна $200\,\text{Дж}$. На какую максимальную высоту поднимется камень? Сопротивлением воздуха пренебречь. Ускорение свободного падения считать равным $10\,\text{м/с}^2$.
}{
\smallskip\noindent\textbf{Решение}\par
Поскольку сопротивлением воздуха можно пренебречь, для камня выполняется закон сохранения полной механической энергии. Вся начальная энергия камня перейдёт в его потенциальную энергию: $E_0=mgh$. Отсюда находим максимальную высоту подъёма камня:
\[h=\frac{E_0}{mg}=\frac{200\,\text{Дж}}{1\,\text{кг}\cdot10\,\text{м/с}^2}=20\,\text{м}.\]
Учитываем, что в точке максимального подъёма скорость камня равна нулю.
Ответ: $20\,\text{м}$.
}

% Source sdamgia-4483: https://phys-ege.sdamgia.ru/problem?id=4483
\BankTask{sdamgia-4483}{
Тело, брошенное вертикально вверх от поверхности Земли, достигло максимальной высоты $20\,\text{м}$. С какой начальной скоростью тело было брошено вверх? Сопротивлением воздуха пренебречь. Ускорение свободного падения считать равным $10\,\text{м/с}^2$.
}{
\smallskip\noindent\textbf{Решение}\par
В момент броска тело имело нулевую потенциальную энергию и не нулевую кинетическую. На высоте максимального подъёма тело не имеет кинетической энергии, поэтому по закону сохранения энергии:
\[\frac{mv^2}{2}=mgh\quad\Longleftrightarrow\quad v=\sqrt{2gh}=\sqrt{2\cdot10\,\text{м/с}^2\cdot20\,\text{м}}=20\,\text{м/с}.\]
Ответ: $20$.
}

% Source sdamgia-9303: https://phys-ege.sdamgia.ru/problem?id=9303
\BankTask{sdamgia-9303}{
У основания гладкой наклонной плоскости шайба массой $10\,\text{г}$ обладает кинетической энергией $0{,}04\,\text{Дж}$. Определите максимальную высоту, на которую шайба может подняться по плоскости относительно основания. Сопротивлением воздуха пренебречь.
}{
\smallskip\noindent\textbf{Решение}\par
На максимальной высоте потенциальная энергия шайбы будет равна её начальной кинетической энергии. Таким образом:
\[mgh=E_{\text{к}}\quad\Longleftrightarrow\quad h=\frac{E_{\text{к}}}{mg}=\frac{0{,}04}{0{,}01\cdot10}=0{,}4\,\text{м}.\]
Ответ: $0{,}4\,\text{м}$.
}

% Source sdamgia-46763: https://phys-ege.sdamgia.ru/problem?id=46763
\BankTask{sdamgia-46763}{
Велосипедист и мотоциклист движутся по горизонтальному мосту. Масса велосипедиста вместе с велосипедом равна $100\,\text{кг}$. Чему равна масса мотоциклиста вместе с мотоциклом, если отношение значений потенциальной энергии мотоциклиста и велосипедиста относительно уровня воды равно $2{,}5$?
}{
\smallskip\noindent\textbf{Решение}\par
Потенциальная энергия тела, поднятого над Землёй, равна $E_{\text{п}}=mgh$. Так как высота, на которой находятся велосипедист и мотоциклист одинаковы, то отношение потенциальных энергий равно
\[\frac{E_{\text{м}}}{E_{\text{в}}}=\frac{m_{\text{м}}}{m_{\text{в}}}=2{,}5,\]
откуда масса мотоциклиста вместе с мотоциклом равна $m_{\text{м}}=2{,}5m_{\text{в}}=2{,}5\cdot100=250\,\text{кг}$.
Ответ: $250\,\text{кг}$.
}

% Source sdamgia-39942: https://phys-ege.sdamgia.ru/problem?id=39942
\BankTask{sdamgia-39942}{
Шарик массой $100\,\text{г}$ падает с некоторой высоты без начальной скорости. К моменту падения на землю потеря полной механической энергии этого шарика из-за действия силы сопротивления воздуха составила $1{,}5\,\text{Дж}$. С какой высоты падал шарик, если кинетическая энергия шарика перед ударом о землю оказалась равна $9{,}5\,\text{Дж}$?
}{
\smallskip\noindent\textbf{Решение}\par
Полная энергия, которой обладал шарик в начале падения, равна его потенциальной энергии и равна $E_{\text{п}}=9{,}5+1{,}5=11\,\text{Дж}$. Тогда высота, с которой падал шарик, равна
\[h=\frac{E_{\text{п}}}{mg}=\frac{11}{0{,}1\cdot10}=11\,\text{м}.\]
Ответ: $11\,\text{м}$.
}

% Source sdamgia-50706: https://phys-ege.sdamgia.ru/problem?id=50706
\BankTask{sdamgia-50706}{
Шайбу с начальной кинетической энергией $E_{\text{к}}=450\,\text{Дж}$ толкают вверх по гладкой горке, она поднимается на максимальную высоту $h=3\,\text{м}$. Определите массу шайбы. Сопротивлением воздуха пренебречь.
}{
\smallskip\noindent\textbf{Решение}\par
При отсутствии силы трения и сопротивления воздуха по закону сохранения энергии $E_{\text{к}1}=E_{\text{п}2}=mgh$, откуда находим массу шайбы:
\[m=\frac{E_{\text{к}1}}{gh}=\frac{450}{10\cdot3}=15\,\text{кг}.\]
Ответ: $15\,\text{кг}$.
}

% Source sdamgia-3598: https://phys-ege.sdamgia.ru/problem?id=3598
\BankTask{sdamgia-3598}{
Небольшое тело массой $0{,}2\,\text{кг}$ бросили вертикально вверх. На рисунке показан график зависимости потенциальной энергии $E_{\text{п}}$ тела от времени $t$ в течение полёта. На какую максимальную высоту поднялось тело?\par\begin{center}
\begin{tikzpicture}[task graph,x=1cm,y=.1cm]
\draw[task grid,xstep=.5,ystep=5] (0,0) grid (5,50);
\draw[task axis] (-.2,0)--(5.15,0) node[below] {$t,\,\text{с}$};
\draw[task axis] (0,-2)--(0,51) node[left] {$E_{\text{п}},\,\text{Дж}$};
\node[below left] at (0,0) {$0$};
\foreach \x in {1,2,3,4} {\fill (\x,0) circle[x radius=1.2pt,y radius=1.2pt]; \node[below] at (\x,0) {$\x$};}
\foreach \y in {10,20,30,40} {\fill (0,\y) circle[x radius=1.2pt,y radius=1.2pt]; \node[left] at (0,\y) {$\y$};}
\draw[task curve,domain=0:4,samples=81] plot (\x,{40-10*(\x-2)^2});
\end{tikzpicture}\end{center}
}{
\smallskip\noindent\textbf{Решение}\par
Потенциальная энергия тела связана с его высотой над поверхностью Земли соотношением $E_{\text{п}}=mgh$. Максимальной высоте полёта соответствует максимум потенциальной энергии. Следовательно, максимальная высота полёта:
\[h_{\max}=\frac{E_{\text{п}\max}}{mg}=\frac{40\,\text{Дж}}{0{,}2\,\text{кг}\cdot10\,\text{м/с}^2}=20\,\text{м}.\]
Ответ: $20\,\text{м}$.
}

% Source sdamgia-3303: https://phys-ege.sdamgia.ru/problem?id=3303
\BankTask{sdamgia-3303}{
Лыжник массой $60\,\text{кг}$ спустился с горы высотой $20\,\text{м}$. Какой была сила сопротивления его движению по горизонтальной лыжне после спуска, если он остановился, проехав $200\,\text{м}$? Считать, что по склону горы он скользил без трения.
}{
\smallskip\noindent\textbf{Решение}\par
Выпишем закон сохранения энергии для незамкнутой системы: $A_{\text{тр}}=E_{\text{пол}2}-E_{\text{пол}1}$, где работа силы трения равна $A_{\text{тр}}=-F_{\text{тр}}s$, полная энергия на высоте горы $E_{\text{пол}1}=E_{\text{п}}=mgh$, полная энергия в момент остановки равна $0$. Следовательно, получаем $F_{\text{тр}}s=mgh$. Таким образом, величина силы трения равна
\[F_{\text{тр}}=\frac{mgh}{s}=\frac{60\,\text{кг}\cdot10\,\text{м/с}^2\cdot20\,\text{м}}{200\,\text{м}}=60\,\text{Н}.\]
Ответ: $60\,\text{Н}$.
}

% Source sdamgia-4207: https://phys-ege.sdamgia.ru/problem?id=4207
\BankTask{sdamgia-4207}{
Маленький шарик начинает падать на горизонтальную поверхность пола с высоты $2\,\text{м}$. Из-за дефектов поверхности пола шарик при ударе о него теряет $20\%$ своей кинетической энергии и отскакивает от пола под углом $60^\circ$ к горизонту. На какую максимальную высоту поднимется шарик после удара о пол?
}{
\smallskip\noindent\textbf{Решение}\par
Будем отсчитывать потенциальную энергию от уровня пола. В начальный момент времени шарик обладает только потенциальной энергией $E=mgh$. Во время падения выполняется закон сохранения энергии, поэтому к моменту удара шарика о пол вся потенциальная энергия переходит в кинетическую. Таким образом, сразу после удара кинетическая энергия шарика становится равной $0{,}8mgh$. После удара шарик летит по параболе, при этом $0{,}8mgh=\frac{mv_0^2}{2}$, откуда начальная скорость шарика равна $v_0=\sqrt{1{,}6gh}$. Проекция начальной скорости на вертикальную ось равна $v_{0y}=v_0\sin60^\circ$. Учитывая, что в верхней точке траектории вертикальная составляющая скорости равна $0$, получаем: $v_y(t)=v_{0y}-gt=0$, следовательно, максимальная высота, на которую поднимется шарик, равна:
\[\begin{aligned}H=\frac{v_{0y}^2}{2g}&=\frac{1{,}6gh\sin^260^\circ}{2g}\\&=0{,}8h\sin^260^\circ=1{,}2\,\text{м}.\end{aligned}\]
Ответ: $1{,}2\,\text{м}$.
}

% Source sdamgia-3410: https://phys-ege.sdamgia.ru/problem?id=3410
\BankTask{sdamgia-3410}{
С балкона высотой $20\,\text{м}$ упал на землю мяч массой $0{,}2\,\text{кг}$. Из-за сопротивления воздуха скорость мяча у земли оказалась на $20\%$ меньше скорости тела, свободно падающего с высоты $20\,\text{м}$. Чему равен импульс мяча в момент удара о землю?
}{
\smallskip\noindent\textbf{Решение}\par
Определим, какую бы скорость приобрёл свободно падающий мяч. По закону сохранения энергии: $\frac{mv^2}{2}=mgh$, откуда
\[v_1=\sqrt{2gh}=\sqrt{2\cdot10\cdot20}=20\,\text{м/с}.\]
Согласно условию, из-за сопротивления воздуха скорость оказалась на $20\%$ меньше: $v_2=0{,}8\cdot20\,\text{м/с}=16\,\text{м/с}$. Таким образом, импульс мяча в момент падения равен $p=mv_2=0{,}2\cdot16=3{,}2\,\text{кг}\cdot\text{м/с}$.
Ответ: $3{,}2\,\text{кг}\cdot\text{м/с}$.
}


% Source B41DCC: https://ege.fipi.ru/bank/questions.php?search=1&proj=BA1F39653304A5B041B656915DC36B38&qid=B41DCC
\BankTask{B41DCC}{
Легковой автомобиль и грузовик движутся по горизонтальному мосту. Масса грузовика $m=4500\,\text{кг}$. Какова масса легкового автомобиля, если отношение потенциальной энергии грузовика к потенциальной энергии легкового автомобиля относительно уровня воды в реке равно $3$? Результат выразите в $\text{кг}$.
}{
\smallskip\noindent\textbf{Решение от GPT}\par
На одном мосту высота над водой одинакова: $E_{\text{п}1}/E_{\text{п}2}=m_1/m_2$. Тогда $m_2=4500/3=1500\,\text{кг}$.
}

% Source 9A9FD1: https://ege.fipi.ru/bank/questions.php?search=1&proj=BA1F39653304A5B041B656915DC36B38&qid=9A9FD1
\BankTask{9A9FD1}{
Спортсмен спускается на парашюте с постоянной скоростью. Как изменяются с течением времени в процессе спуска импульс спортсмена и его потенциальная энергия?
\begin{enumerate}
\item увеличивается
\item уменьшается
\item не изменяется
\end{enumerate}
\TaskChangeTable{Импульс}{Потенциальная энергия}
}{
\smallskip\noindent\textbf{Решение от GPT}\par
При постоянной скорости $\vec p=m\vec v$ не меняется. Высота уменьшается, поэтому $E_{\text{п}}=mgh$ уменьшается. Ответ: $32$.
}

\end{document}
